\anexo{Entrevistas: guía y perfiles}\label{annex:entrevistas}

Este anexo reúne la guía de entrevista y los perfiles de los tres participantes. Las transcripciones completas de las conversaciones, mantenidas el 22 de octubre de 2026 por videollamada, figuran en los Anexos~\ref{annex:entrevista-tadeo}, \ref{annex:entrevista-fabricio} y~\ref{annex:entrevista-octavio}. Cada registro consigna fecha, hora de inicio, entrevistador y duración. Las transcripciones son automáticas y se reproducen con una edición mínima: se corrigen entre corchetes los errores evidentes del reconocimiento de voz (por ejemplo, el broker Balanz transcripto como \enquote{Balance}) y, en la tercera, se suprimen las muletillas de inicio de frase sin alterar el contenido. Al inicio de cada conversación el entrevistador informa que las respuestas se usan para el Proyecto, lo que consta en la propia transcripción.

\section*{Guía utilizada}
La guía organiza la conversación en cinco ejes y se transcribe tal como se aplicó.
\begin{enumerate}
\item ¿Cómo es hoy tu manera de seguir y analizar tus inversiones? ¿Qué herramientas, apps o planillas usás y cómo las combinás?
\item Cuando tenés activos de distinto tipo (por ejemplo, acciones, ETFs y criptomonedas), ¿tus inversiones están todas en un mismo lugar o repartidas en varias apps o cuentas?
\item La última vez que tomaste una decisión de comprar, vender o reacomodar tu cartera, ¿en qué información te basaste y qué te hubiera gustado tener y no tenías?
\item ¿Alguna vez una app o plataforma te dio una alerta, sugerencia o recomendación sobre tus inversiones? ¿Entendiste por qué te la daba y qué tanto confiaste en ella?
\item Si pudieras pedirle cualquier cosa a una herramienta ideal para analizar tu portafolio, ¿qué te gustaría que hiciera por vos? ¿Y qué cosas no querrías que hiciera?
\end{enumerate}

\section*{Perfiles de los participantes}
Los tres participantes son inversores minoristas del contexto argentino, ajenos al equipo y a la tutoría, seleccionados por conveniencia entre contactos de los autores con experiencia de inversión en más de una plataforma. La selección busca cubrir el segmento activo de la encuesta, que es el que más necesita consolidar información.
\begin{itemize}
\item \textbf{Tadeo Pinque.} Opera acciones y ETFs desde la aplicación del intermediario; recurre a otras plataformas cuando un instrumento no está disponible desde la Argentina. Revisa resultados a fin de mes con una herramienta de IA. Entrevistado por Mauricio Huentelaf.
\item \textbf{Fabricio Sabattoli.} Mantiene activos tradicionales y renta fija en Balanz y criptoactivos en Binance; consolida de forma manual en Yahoo Finance, donde no puede cargar bonos, letras ni fondos locales. Se informa por X y Yahoo Finance y usa Claude como apoyo. Entrevistado por Máximo Rosso.
\item \textbf{Octavio Antunes.} Reparte acciones y ETFs en Interactive Brokers, criptoactivos en Binance y otras inversiones en un broker local; lleva un registro en Excel y consulta Yahoo Finance, Google Finance y TradingView. Entrevistado por Mauricio Huentelaf.
\end{itemize}

\anexo{Transcripción de la entrevista a Tadeo Pinque}\label{annex:entrevista-tadeo}

Fecha: 22 de octubre de 2026, 19:36. Entrevistador: Mauricio Huentelaf. Duración: 3 minutos 31 segundos.
\par\noindent\textbf{[00:01] Mauricio Huentelaf:} Bueno Tadeo, buenas, te voy a hacer unas preguntas ahora para nuestro proyecto de tesis. La primera sería ?`Cómo es tu manera de seguir y analizar tus inversiones y qué herramientas, apps o planillas usas y cómo las combinas?
\par\noindent\textbf{[00:22] Tadeo Pinque:} Bueno, en primer lugar, para seguir mis inversiones, utilizo la app en la cual compro las acciones y los ETFs y, Y a su vez a fin de mes utilizo una inteligencia artificial para de esta manera analizar los resultados y validar en que podría mejorar.
\par\noindent\textbf{[00:42] Mauricio Huentelaf:} Bien, perfecto. Cuando tenías activos de distinto tipo, por ejemplo acciones, ETF y criptomonedas, ?`tus inversiones están todas en un mismo lugar o repartidas en varias apps y cuentas?
\par\noindent\textbf{[00:57] Tadeo Pinque:} En realidad, usualmente trato de utilizar siempre una misma plataforma de inversión, aunque bueno, a veces termino usando distintas plataformas ya que hay acciones que no se pueden comprar desde Argentina o incluso criptomonedas que no están disponibles en ciertos brokers.
\par\noindent\textbf{[01:18] Mauricio Huentelaf:} Ok, buenísimo. La última vez que tomaste una decisión de comprar, vender o reacomodar tu cartera, ?`en qué información te basaste y, Y qué te hubiera gustado tener y no tenías?
\par\noindent\textbf{[01:31] Tadeo Pinque:} Bueno, la última vez que compré recuerdo que fue una acción de semiconductores, ya que creo que tienen un gran impacto y potencial de crecimiento.
\par\noindent\textbf{[01:44] Tadeo Pinque:} La compré aprovechando una caída que tuvo recientemente y me basé en distintos análisis de mercado y también en chats con inteligencia artificial. Tratando de entender el rumbo de la empresa y su potencial. Me hubiera gustado también que la misma plataforma me diera cierta información sobre estas acciones para poder entenderlas mejor antes de tomar una decisión.
\par\noindent\textbf{[02:10] Mauricio Huentelaf:} Alguna vez una aplicación o plataforma te dio una alerta, sugerencia o recomendación sobre tus inversiones?
\par\noindent\textbf{[02:23] Tadeo Pinque:} Hay varias plataformas que te muestran el análisis o la opinión de sus inversores. Esto lo tomo como un parámetro, pero igualmente me gusta hacer mis propias investigaciones para validar que esas afirmaciones sean correctas.
\par\noindent\textbf{[02:37] Mauricio Huentelaf:} Bien, excelente. Bueno, y por último, si pudieras pedirle a cualquier persona una herramienta ideal para analizar tu portafolio, ?`qué te gustaría que hiciera por vos y qué cosas no querrías que hiciera?
\par\noindent\textbf{[02:52] Tadeo Pinque:} Me gustaría que me recomiende cómo reorganizar mi portafolio para gestionar distintos estilos de inversión. Por ejemplo, si soy un inversor con un poco menos de aversión al riesgo y tengo muchas acciones relacionadas con IA que me pueda decir algo, como por ejemplo, para poder compensar un poco tu portafolio podrías incorporar acciones de otro estilo, como por ejemplo, Berkshire Hathaway o alguna otra demás por el estilo. La idea sería tener un, Una cartera más equilibrada y en caso de que el mercado castigue fuerte no tenga mucha pérdida.
\par\noindent\textbf{[03:31] Mauricio Huentelaf:} Excelente Tadeo, bueno muchísimas gracias por responder a las preguntas.

\anexo{Transcripción de la entrevista a Fabricio Sabattoli}\label{annex:entrevista-fabricio}

Fecha: 22 de octubre de 2026, 19:45. Entrevistador: Máximo Rosso. Duración: 3 minutos 35 segundos.
\par\noindent\textbf{[00:00] Máximo Rosso:} Hola Fabricio, ?`Cómo estás? Bueno, empezamos con la entrevista. La primera pregunta que te voy a hacer es ?`Cómo es tu manera de seguir y analizar tus inversiones? Qué herramientas, aplicaciones o planillas usas y cómo las combinas?
\par\noindent\textbf{[00:17] Fabricio Sabattoli:} Bien, principalmente uso Balanz [transcripto como \enquote{Balance}], ya que tengo prácticamente todas mis inversiones ahí. La otra parte lo tengo en Binance, algo que hago, Para seguir gran parte de mis inversiones es subir todo a Yahoo Finance, donde puedo tener todo en uno o casi todo, ya que instrumentos financieros como bonos, letras o fondos que tengo en Balanz no se encuentran en Yahoo Finance. Pero uso este último más que nada para leer noticias y mantenerme al día.
\par\noindent\textbf{[00:52] Máximo Rosso:} Perfecto, te voy a hacer la segunda pregunta. Cuando tenes activos de distintos tipos, Por ejemplo, acciones, ETFs y criptomonedas. Tus inversiones están todas en un mismo lugar o repartidas en varias apps o cuentas?
\par\noindent\textbf{[01:06] Fabricio Sabattoli:} No, generalmente las tengo repartidas entre Binance y Balanz [transcripto como \enquote{Balance}].
\par\noindent\textbf{[01:13] Máximo Rosso:} Ok, perfecto. La última vez que tomaste una decisión de comprar, vender o reacomodar tu cartera. En qué información te basaste y qué te hubiera gustado tener y no tenías?
\par\noindent\textbf{[01:26] Fabricio Sabattoli:} Bueno, en este caso yo me baso en información que recopilo en X y Yahoo Finance principalmente, pero no siento que me falte información y también uso herramientas como por ejemplo Claude [transcripto como \enquote{Cloud}].
\par\noindent\textbf{[01:44] Máximo Rosso:} Exacto, Fabricio. Alguna vez alguna aplicación o plataforma te dio una alerta, sugerencia o recomendación sobre tus inversiones? La entendiste por qué te la daba y qué tanto confiaste en ella?
\par\noindent\textbf{[01:57] Fabricio Sabattoli:} Nunca una recomendación de inversión como tal, pero sí hay muchos post en X, los cuales indirectamente se inclinan hacia una recomendación y ahí depende de uno si tomarla o rechazarla. Si dicho post se alinea o refuerza la posición de uno en una empresa o activo, suelo tomarlo como recomendación de inversión. También hay que entender que cuando uno ve un post o noticia sobre una recomendación de inversión, la persona o empresa detrás ya tomó una posición previa, por ejemplo, cuando un banco cambia la recomendación de una empresa a Underweight o Underperform, lo más probable es que dicho banco ya haya vendido o reducido su posición en dicha empresa.
\par\noindent\textbf{[02:53] Fabricio Sabattoli:} Y el efecto ya esté desencadenado en el precio, básicamente.
\par\noindent\textbf{[02:59] Máximo Rosso:} Perfecto, Fabricio. Por último, si pudieras pedirle cualquier cosa a una herramienta ideal para analizar tu portafolio, ?`qué te gustaría que hiciera por vos y qué cosas no querrías que haga?
\par\noindent\textbf{[03:12] Fabricio Sabattoli:} Bueno, me gustaría poder pedirle el sentimiento sobre cierto activo para así poder ver ambos lados de la moneda, Y que la decisión de inversión sea más robusta.
\par\noindent\textbf{[03:28] Fabricio Sabattoli:} Y por otro lado, no se me ocurre algo que no me gustaría que hiciera básicamente.
\par\noindent\textbf{[03:35] Máximo Rosso:} Ok, perfecto.

\anexo{Transcripción de la entrevista a Octavio Antunes}\label{annex:entrevista-octavio}

Fecha: 22 de octubre de 2026, 19:52. Entrevistador: Mauricio Huentelaf. Duración: 5 minutos 57 segundos.
\par\noindent\textbf{[00:01] Mauricio Huentelaf:} Buenas Octavio, ?`cómo andas? Ahora te voy a hacer unas preguntas para nuestro proyecto de tesis. Para arrancar, ?`cómo es tu manera de seguir y analizar tus inversiones? Qué herramientas, aplicaciones o planillas usas y cómo las combinas?
\par\noindent\textbf{[00:24] Octavio Antunes:} En primer lugar, actualmente hago el seguimiento de mis inversiones de manera, bastante manual. Utilizo aplicaciones de los brokers en todos mis activos y también herramientas como por ejemplo Yahoo Finance y Google Finance o TradingView para consultar precios gráficos y noticias del mercado. Además, en algunas ocasiones utilizo un Excel como para llevar un registro de cuánto invertí y cómo evolucionó mi cartera. Mi principal problema es que la información está bastante dispersa y tengo que estar.
\par\noindent\textbf{[00:57] Octavio Antunes:} Al tanto y entrar a distintas plataformas para entender como están funcionando mis inversiones ademas siento que las aplicaciones me muestran muchos números pero no necesariamente me ayudan a interpretarlos o a tomar mejores decisiones.
\par\noindent\textbf{[01:12] Mauricio Huentelaf:} Ok perfecto cuando tienes activos de distinto tipo por ejemplo acciones, ETF o criptomonedas tus inversiones estan todas en un mismo lugar o repartidas en varias aplicaciones, cuentas, etc.
\par\noindent\textbf{[01:28] Octavio Antunes:} Mis inversiones están repartidas en varias plataformas. Por ejemplo, yo utilizo Interactive Brokers para acciones y ETFs y Binance para criptomonedas. También tengo algunas inversiones a través de un broker local. Esto hace que sea difícil tener una visión completa de mi patrimonio invertido. Cada aplicación me muestra su propio rendimiento, pero ninguna me permite entender fácilmente cómo está distribuida toda mi cartera. Para calcular qué porcentaje tengo en acciones, ETFs o criptomonedas, generalmente tengo que hacerlo manualmente.
\par\noindent\textbf{[02:02] Octavio Antunes:} Pero me gustaría poder visualizar todo esto dentro de un mismo lugar, sin la necesidad de transferir mis activos de las plataformas en las que yo utilizo.
\par\noindent\textbf{[02:12] Mauricio Huentelaf:} Perfecto. La última vez que tomaste una decisión de comprar, vender o reacomodar tu cartera, ?`en qué información te basaste y qué te hubiera gustado tener? No tenías?
\par\noindent\textbf{[02:27] Octavio Antunes:} La última vez que tomé una decisión de inversión fue para aumentar una exposición al mercado americano a través de un ETF del S\&P 500. Antes de comprarlo, analicé su rendimiento histórico, las condiciones generales del mercado, algunas noticias económicas y la evolución de las tasas de interés. También comparé esta alternativa con invertir directamente en acciones de empresas tecnológicas. De igual manera, considero que me faltó información sobre cómo esta inversión afectaría el riesgo total de mi cartera, pero yo sabía que estaba comprando, pero no tenía una herramienta que me permitiera analizar fácilmente si estaba demasiado concentrado en determinados sectores o mercados.
\par\noindent\textbf{[03:12] Octavio Antunes:} Me hubiera resultado útil contar con una comparación de escenarios y alguna explicación sobre cómo esa compra que hice modificaría la diversificación. Y el riesgo de mis inversiones.
\par\noindent\textbf{[03:25] Mauricio Huentelaf:} Bien. Alguna vez una aplicación o plataforma te dio una alerta, sugerencia o recomendación sobre tus inversiones? Entendiste qué te la daba y qué tanto confiaste en ella?
\par\noindent\textbf{[03:44] Octavio Antunes:} Sí, recibí alertas relacionadas con movimientos importantes de precios y algunas sugerencias sobre activos que estaban teniendo un buen rendimiento. Mi problema es que muchas veces esas notificaciones son demasiado generales. Por ejemplo, una aplicación puede avisarme con una acción cayó un 5\% pero no me explica si esa caída [transcripto como \enquote{calidad}] está relacionada con resultados financieros de la empresa, una noticia política o un movimiento del mercado general. Por este motivo, no suelo tomar decisiones únicamente para partir de esas alertas, lo utilizo como punto de partida para investigar por mi propia cuenta.
\par\noindent\textbf{[04:19] Octavio Antunes:} Confiaría mucho más en una recomendación si pudiera entender claramente qué información se utilizó para generarlo, qué riesgo considera y cómo se relaciona con mi cartera particular.
\par\noindent\textbf{[04:31] Mauricio Huentelaf:} Bien, y por último, si pudieras pedirle cualquier cosa o una herramienta ideal para analizar tu portafolio, ?`qué te gustaría que hiciera por vos y qué cosas no querrías que hiciera?
\par\noindent\textbf{[04:45] Octavio Antunes:} Lo principal, Quería que pudiera integrar todas mis inversiones en una sola plataforma, independientemente de donde tenga depositados mis activos. Me gustaría que no solamente mostrara precios y rendimientos, sino que también analizara la composición de mi cartera, identificara riesgos de concentración y me ayudara a entender cómo determinados acontecimientos económicos o políticos podrían afectar mis inversiones. Por otra parte, valoraría que me notificara cuando ocurriera algo realmente relevante para los activos que tengo, en lugar de enviarme alertas constantemente sobre movimientos que quizás no requieran ninguna acción.
\par\noindent\textbf{[05:27] Octavio Antunes:} Por otro lado, no me gustaría que la herramienta comprara o vendiera activos automáticamente sin mi propia autorización, y tampoco confiaría en un sistema que prometiera rendimientos garantizados o que me presentara recomendaciones sin explicar sus, sus fundamentos. En definitiva, buscaría una herramienta que me ayudara a tomar decisiones más informadas, pero que a su vez no pretendiera reemplazar mi propio criterio como inversor.
\par\noindent\textbf{[05:57] Mauricio Huentelaf:} Bien, buenísimo. Muchas gracias, Octavio.
